\section{曲面上的多边形}\label{sec:13}\label{section:13}
\subsection{自动正则性}\label{sec:13a}
设 $S$ 为带条带状端的连通 Riemann 曲面。考虑 $S$ 上所有在某个紧子集之外与给定 $j_0=I_S$ 一致的复结构 $j$ 所成的空间 $\widetilde{\RG}$。乘积 $\widetilde{\mathcal S}=\widetilde{\RG}\times S\to\widetilde{\RG}$ 的纤维上有自然的复结构，因此成为带条带状端的 Riemann 曲面族。此外，该族上有群 $\widetilde{\mathcal G}$ 的自然作用，该群由 $S$ 的如下微分同胚组成：在每个条带状端的无穷远附近，它们限制为 $s$ 方向的平移。

为 $S$，因而也为 $\widetilde{\mathcal S}$，赋予 Lagrangian 标记，使每个无穷远点 $\zeta$ 都满足 $L_{\zeta,\allowbreak 0}\pitchfork L_{\zeta,\allowbreak 1}$。选择平凡扰动数据 $K=0$、$J=I_M$，即殆复结构的常值族。取伪全纯映射 $u\in\M_S(\{y_\zeta\})$，将 $(j_0,\allowbreak u)$ 视为相应参数化模空间 $\M_{\widetilde{\mathcal S}}(\{y_\zeta\})$ 的元素。$\widetilde{\mathcal G}$ 的无穷小作用定义从 Lie 代数 $L\widetilde{\mathcal G}$ 到 Zariski 切空间 $T^{\mathrm{Zar}}\M_{\widetilde{\mathcal S}}(\{y_\zeta\})_{(j_0,\allowbreak u)}$ 的线性映射 $d^0$；使用“Zariski 切空间”这一说法，是因为尚未作正则性假设。另一方面，依定义该 Zariski 切空间是扩展线性化算子 $d^1=D_{\widetilde{\mathcal S},\allowbreak j_0,\allowbreak u}$ 的核。因此可将两个算子视为如下复形的边界映射：
\begin{equation}\label{eq:13.1}
0\longrightarrow L\widetilde{\mathcal G}\xrightarrow{d^0}T\widetilde{\RG}_{j_0}\times(T\mathcal B_S)_u
\xrightarrow{d^1}(\mathcal E_S)_u\longrightarrow0.
\end{equation}
到目前为止，对无穷维空间 $\widetilde{\RG}$、$\widetilde{\mathcal G}$ 的处理相当非正式；现在至少在切空间层面，希望通过选择合适的 Hilbert 空间完备化使其精确。按先前定义，$\widetilde{\mathcal G}$ 的 Lie 代数由 $S$ 上与边界相切的光滑向量场 $Z$ 组成，它们在每个条带状端的无穷远附近满足 $\epsilon_\zeta^*Z=c_\zeta\partial_s$，其中 $c_\zeta\in\R$ 为常数。现在改取 $(k,\allowbreak p)=(2,\allowbreak 2)$、权重为充分小 $\epsilon>0$ 的加权 Sobolev 完备化，即\eref{9.2}记号中的 $W^{2,\allowbreak 2;\epsilon}(S,\allowbreak TS,\allowbreak T(\partial S))^+$。同样，朴素地说，$\widetilde{\RG}$ 的切空间是与复结构反交换的 $TS$ 的紧支自同态空间。将这些自同态解释为复张量积线丛 $\Omega_S^{0,\allowbreak 1}\otimes TS$ 的光滑截面，并完备化为
\SourcePage{199}
$W^{1,\allowbreak 2;\epsilon}(\Omega_S^{0,\allowbreak 1}\otimes TS)$。最后，对作为 $d^1$ 第二个分量的线性化算子 $D_{S,\allowbreak u}$，使用相同的 Sobolev 拓扑。\eref{13.1}的完备化于是为
\begin{equation}\label{eq:13.2}
\begin{gathered}
W^{2,2;\epsilon}(S,TS,T(\partial S))^+\\
\big\downarrow d^0\\
W^{1,2;\epsilon}(S,\Omega_S^{0,1}\otimes TS)\oplus W^{1,2;\epsilon}(S,E,F)\\
\big\downarrow d^1\\
L^{2;\epsilon}(S,\Omega_S^{0,1}\otimes E),
\end{gathered}
\end{equation}
其中 $E=u^*TM$，$F\subset E|\partial S$ 由相关 Lagrangian 子流形的切空间给出，且
\[
d^0(Z)=(2I_S\circ\bar\partial Z,Du\circ Z),\qquad
d^1(Y,X)=\tfrac12I_M\circ Du\circ Y+D_{S,u}(X).
\]
由于 Sobolev 空间的选择是为当前问题临时作出的，首先应验证这些算子确实有界，不过这很直接。事实上，还能得到更多结论：$TS$ 上边界条件为 $T(\partial S)$ 的标准 Cauchy–Riemann 算子 $\bar\partial$ 容许，故在加权 Sobolev 空间上成为 Fredholm 算子，见\secref{9b}。$D_{S,\allowbreak u}$ 本身是非退化 $\bar\partial_\nabla$ 算子，因此也为 Fredholm 算子。其余分量为零阶项，因而紧。结论是\eref{13.2}为 Fredholm 复形。

选择曲面族 $\mathcal S\to\RG$，这次底空间有限维，使它成为\secref{9b}意义下 $S=\mathcal S_{r_0}$ 的无穷小极小通用变形。令 $\mathcal G$ 为 $S$ 的全纯自同构群。
\begin{lemma}\label{st:13.1}
\eref{13.2}的上同调群具有如下解释：$H^0,\allowbreak H^1$ 分别为映射 $L\mathcal G\to\ker(D_{\mathcal S,\allowbreak r_0,\allowbreak u})$ 的核与余核；$H^2=\operatorname{coker}(D_{\mathcal S,\allowbreak r_0,\allowbreak u})$。
\end{lemma}
\begin{proof}
非正式地，想法是 $\RG$ 为 $\widetilde{\mathcal G}$ 在 $\widetilde{\RG}$ 上作用的局部切片。为使论证严格，选择在条带状端上恒定的 $\mathcal S$ 的局部平凡化。这给出 $\RG$ 到 $\widetilde{\RG}$ 的局部嵌入，因而给出
\SourcePage{200}
切空间之间的映射。于是有各列正合的交换图
\SourceEquationDiagram{diagram-13-3}{0.36969191}{13.3}
依定义，$T\RG_{r_0}$ 同构地投影到\eref{9.3}的余核，因此\eref{13.3}的底行是满足 $H^0=L\mathcal G$、$H^1=0$ 的复形。上同调长正合列给出所需结论。
\end{proof}
\begin{lemma}\label{st:13.2}
若 $M$ 为二维且 $u$ 非常值，则对任意无穷小极小通用族 $\mathcal S\to\RG$，$D_{\mathcal S,\allowbreak r_0,\allowbreak u}$ 均为满射。
\end{lemma}
\begin{proof}
$u$ 为 Riemann 曲面之间的非常值全纯映射，所以除有限多个点外，$Du:TS_z\to TM_{u(z)}$ 都是同构。因此 $d^1$ 的第一个分量具有稠密像。由于第二个分量为 Fredholm 算子，便知 $d^1$ 满射。再运用\stref{引理}{13.1}最后一条即可。
\end{proof}
称\stref{引理}{13.2}为“自动正则性”结论，因为它对任意 $I_M$ 成立，且不需要扰动伪全纯映射方程。\znote{192}{自动正则性的关键在 $\dim M=2$：此时 $S$ 与 $M$ 的切空间都是复一维，$u$ 非常值全纯就使 $Du$ 除有限点外为同构，于是 $d^1$ 的第一个分量已提供足够多的方向，第二个 Fredholm 分量再补足余下部分，满射性无需额外假设。若 $M$ 维数更高，$Du$ 只是单射，$d^1$ 的像一般并不稠密，必须靠扰动才能保证正则。这也解释了为什么组合 Floer 理论只在曲面情形如此干净。[来源：Seidel, Fukaya categories and Picard–Lefschetz theory, §13a 与引理 13.2。]}
\begin{note}\label{st:13.3}
偏好代数几何的读者也许会觉得下面的层论解释有帮助。为简单起见，转到类似的情形：全纯映射 $u:S\to M$，其中 $S$ 是闭 Riemann 曲面，$M$ 是复流形。此时\eref{13.2}变为
\[
\begin{gathered}
W^{2,2}(S,TS)\longrightarrow
W^{1,2}(S,\Omega_S^{0,1}\otimes TS)\oplus W^{1,2}(S,u^*TM)\\
\longrightarrow L^2(S,\Omega_S^{0,1}\otimes u^*TM).
\end{gathered}
\]
它实际上计算凝聚层的两项复形 $\{Du:TS\to u^*TM\}$ 的超上同调。若 $u$ 非常值，则 $Du$ 作为层映射为单射，所以实际上计算的是商层 $u^*TM/TS$ 的上同调，再平移一阶。最后，若 $M$ 是曲面，则该商为挠层，因此一次上同调消失。
\end{note}

\subsection{组合 Floer 上同调}\label{sec:13b}
设 $M$ 为二维带角正合辛流形。$M$ 中连通正合 Lagrangian 子流形是
\SourcePage{201}
满足 $\int_L\theta=0$ 的简单闭曲线 $S^1\cong L\subset\operatorname{int}(M)$，其中 $\theta\in\Omega^1(M)$ 为辛形式 $\omega$ 的指定原始形式。若曲线有多个连通分支，则每个分支上的积分均须为零。注意，由 Stokes 定理，这排除了零同调的曲线。固定一族相互处于一般位置的正合 Lagrangian 子流形 $\symbfit{L}=(L_1,\allowbreak \ldots,\allowbreak L_m)$，即任意两条横截相交，且无三重交点。定义\sr{sec:9}{第~9~节}初步版本的、不分次也不带符号的 Fukaya 范畴 $\F(M)^{\mathrm{pr}}$ 时，可以将 Floer 数据
\begin{equation}\label{eq:13.4}
(H_{L_{i_0},L_{i_1}},J_{L_{i_0},L_{i_1}}),\qquad i_0<i_1,
\end{equation}
取为平凡数据 $H=0,\allowbreak J=I_M$，扰动数据
\begin{equation}\label{eq:13.5}
(K_{L_{i_0},\ldots,L_{i_d}},J_{L_{i_0},\ldots,L_{i_d}}),\qquad i_0<\cdots<i_d,
\end{equation}
也可如此选择。验证这一做法可行涉及两方面。第一方面由\stref{引理}{13.2}解决，它保证对应于\eref{13.4}和\eref{13.5}的 Floer 轨线模空间及全纯多边形模空间分别正则。这些模空间没有常值映射，因为边界上的 Lagrangian 标记 $L_{i_k}$ 各不相同且处于一般位置。第二方面是\eref{13.5}与万有扰动数据选择的一致性条件相容，见\secref{9i}。这是以下初等观察的推论：取稳定 $d$ 叶树 $T$，并对每个顶点取带 Lagrangian 标记和条带状端的 $|v|$ 点圆盘 $S_v$。假设标记相容，使各 $S_v$ 可粘合为带 Lagrangian 标记的 $(d+1)$ 点圆盘 $S$。若 $S$ 的标记为 $(L_{i_0},\allowbreak \ldots,\allowbreak L_{i_d})$，且 $i_0<\cdots<i_d$，则每个组成部分 $S_v$ 的标记均为其有序子集，因而具有同样形式。

接着，假定\eref{13.4}与\eref{13.5}均取平凡数据。断言：\secref{5n}意义下的有向 $\Ai$ 子范畴 $\F(\symbfit{L})^{\mathrm{pr},\allowbreak \to}\subset\F(M)^{\mathrm{pr}}$ 与复结构 $I_M$ 无关。换言之，固定 $i_0<\cdots<i_d$ 及交点 $y_0\in L_{i_0}\cap L_{i_d}$、$y_k\in L_{i_{k-1}}\cap L_{i_k}$（$1\leq k\leq d$），则 $\M^{d+1}(y_0,\allowbreak \ldots,\allowbreak y_d)$ 中孤立点的个数在 $I_M$ 的同伦下不变。这由模空间的直接配边论证证明，也可视为\secref{10e}讨论的平凡情形：由\stref{引理}{13.2}，参数化模空间 $\M^{d+1}(y_0,\allowbreak \ldots,\allowbreak y_d)^{[0;1]}$ 光滑纤维化于 $[0;1]$ 上，故不含孤立点，因此\stref{命题}{10.13}的 $\Ai$ 函子为恒等函子。

既已建立 $\F(\symbfit{L})^{\mathrm{pr},\allowbreak \to}$ 原则上具有拓扑性质，下一步是得到显式公式。这方面已有充分研究，见本章开头的参考文献，因此这里只作简述。令 $\Delta^{d+1}\subset\R^2$ 为有定向的凸 $(d+1)$ 边形；$d=1$ 时为二边形，或半月形。其角点和边的编号与 $(d+1)$ 点圆盘相同。对前述点 $y_0,\allowbreak \ldots,\allowbreak y_d$，
\SourcePage{202}
考虑保持正定向的光滑浸入 $\Delta^{d+1}\to M$ 所成的集合 $\operatorname{Imm}^{d+1}(y_0,\allowbreak \ldots,\allowbreak y_d)$，要求第 $k$ 个角点映到 $y_k$，第 $k$ 条边映到 $L_{i_k}$，见\sr{fig:13.1}{图~13.1}。坚持要求映射直到边界都是浸入，
\SourceFigure{figure-13-1}{0.40084572}{13.1}
即每个角点附近的像是顶角小于 $\pi$ 的楔形区域。固定各角点的 $\Delta^{d+1}$ 微分同胚群作用于 $\operatorname{Imm}^{d+1}(y_0,\allowbreak \ldots,\allowbreak y_d)$，其商可与 $\M^{d+1}(y_0,\allowbreak \ldots,\allowbreak y_d)$ 的零维部分识别。证明如下。为方便，假定选择 $I_M$，使每个两曲线交点附近都有局部全纯坐标，在其中两曲线成为 $\CC$ 中的实轴和虚轴。这没有问题，因为已知 $\F(\symbfit{L})^{\mathrm{pr},\allowbreak \to}$ 与 $I_M$ 无关；也可处理任意 $I_M$，代价是多用一些分析，如\bcite{114}所示。指标论证\bcite{43}证明，每个孤立 $u\in\M^{d+1}(y_0,\allowbreak \ldots,\allowbreak y_d)$ 都是浸入，且其像在无穷远点附近呈楔形。再稍加考察 $u$ 在条带状端上的收敛速率，便知可将其紧化为浸入 $\tilde u:\Delta^{d+1}\to M$。反方向从浸入出发构造全纯映射，依赖以下事实：在 $\Delta^{d+1}$ 上给定任意复结构后，去掉角点，剩下的就是 $(d+1)$ 个边界穿孔的圆盘；这很容易归结为闭圆盘的单值化定理。\znote{194}{这一步把分析问题化成组合问题：$\Delta^{d+1}$ 上的任意复结构在去掉角点后给出边界穿孔的圆盘，而单值化定理保证它与标准圆盘共形等价，边界点按顺序对应。于是多边形的复结构不带来额外模参数，零维模空间的点与浸入多边形及其角点重参数化一一对应，$\mu^d$ 因而可以只靠数多边形来计算。13c 能用纯组合数据定义 $\tilde\mu^d$，依据正在于此。[来源：Seidel, Fukaya categories and Picard–Lefschetz theory, §13b（原文所引单值化定理）。]}
\begin{note}\label{st:13.4}
自然会问，由一般位置的一族 $\symbfit{L}$ 组成的 $\F(M)^{\mathrm{pr}}$ 的满 $\Ai$ 子范畴，是否也有组合描述。看来很可能有，但答案应比这里给出的复杂得多，
\SourcePage{203}
因为交点处的常值全纯多边形也有贡献，而此时\stref{引理}{13.2}失效：其模空间的维数过大，虽光滑，却不正则。\znote{195}{常值映射是唯一的例外：$u\equiv y$ 时 $Du=0$，引理 13.2 证明中“$Du$ 除有限点外为同构”完全失效，$D_{S,u}$ 的满射性没有保障。几何上这些点对应退化多边形：所有角点都落在同一交点附近，多边形缩成一点。它们在一般位置下仍对 $\mu^d$ 有贡献，所以 $\F(M)^{\mathrm{pr}}$ 的满 $\Ai$ 子范畴不能简单取组合模型，这正是 13b 只处理有向子范畴、而不直接采用组合定义的原因。[来源：Seidel, Fukaya categories and Picard–Lefschetz theory, §13b，Remark 13.4 与引理 13.2。]}
\end{note}
